\section{A Broader Shift in AI Systems}

Fusion is part of a broader movement from monolithic model invocation toward
systems that actively shape inference. Routing and cascade systems demonstrate
that quality and cost can be optimized jointly
\citep{chen2023frugalgpt,ong2024routellm,aggarwal2024automix}. Ensembling and
Mixture-of-Agents systems show that multiple candidates can improve a final
answer when their contributions are complementary
\citep{jiang2023llmblender,wang2024moa,lu2026complementary}. Learned
orchestrators increasingly treat worker selection, roles, and topology as
policy decisions
\citep{zhang2025routerr1,yue2025masrouter,nielsen2026conductor,xu2026trinity,fugu2026}.
Long-horizon agents extend the problem from requests to evolving trajectories
\citep{zhang2026mtrouter,koh2026macu}.

Fusion shares the conviction that inference should be adaptive and learnable,
but chooses a broader control object. A router asks which model should answer.
An ensemble asks which answers should be combined. An orchestrator often asks
which worker graph should execute. Fusion asks what decision-relevant evidence
the current task needs next, which capability can provide it, and whether the
system is ready to commit.

That distinction connects several previously separate concerns. Model choice,
role, fan-out, tool use, verification, stopping, cost, and workflow phase all
become aspects of one allocation decision. The conceptual roots include active
information acquisition and value-of-information reasoning
\citep{houlsby2011bald,golovin2011adaptive}; the system contribution is to make
those ideas practical under natural-language ambiguity, correlated model
errors, heterogeneous operating costs, and changing external state.

The important transition is not from one model to many models. It is from
static consumption of intelligence to adaptive allocation of intelligence.
